%% toc: P1 8 8.1,8.2,8.5
\section{Proofs for the readings of the Part 9 program constructors and diamonds, the autoepistemic operator and the Part 19 denial}
\label{sec:annex-p01-w02c-p01-017-readings}

\begin{annexitem}{annex:p01-w02c-p01-017-pdl-bracketing}{thm:p01-w02c-p01-017-pdl-bracketing}{Bracketing does not affect the relation of an iterated Part~9 composition or choice}
\begin{proof}
Fix the model \(M\).  The only facts about programs that the argument uses
are the two clauses in the statement: for all programs \(\pi,\rho\) and
states \(w,v\), \((w,v)\in R_{\pi;\rho}\) iff \((w,u)\in R_\pi\) and
\((u,v)\in R_\rho\) for some state \(u\), and
\(R_{\pi\cup\rho}=R_\pi\cup R_\rho\).  \Cref{def:p09-pdl} gives them for
arbitrary programs, including programs that contain tests and iterations,
since composition and union denote relational composition (left program
first) and union.

A bracketing of \(n\geq1\) operands in a fixed order is a single operand
when \(n=1\), and otherwise joins a bracketing of \(\pi_1,\ldots,\pi_k\)
to a bracketing of \(\pi_{k+1},\ldots,\pi_n\) for some \(k\) with
\(1\leq k<n\).  We prove both items by induction on \(n\), for all states
\(w,v\) and all programs \(\pi_1,\ldots,\pi_n\).

\emph{Case \(n=1\).}  The only bracketing of either kind is \(\pi_1\).  For
item~1, states \(u_0,u_1\) with \(u_0=w\), \(u_1=v\) and
\((u_0,u_1)\in R_{\pi_1}\) exist iff \((w,v)\in R_{\pi_1}\).  Item~2 reads
\(R_{\pi_1}=R_{\pi_1}\).

\emph{Inductive step for item~1.}  Let \(n>1\) and
\(\gamma=\gamma_1;\gamma_2\), where \(\gamma_1\) brackets
\(\pi_1,\ldots,\pi_k\) and \(\gamma_2\) brackets
\(\pi_{k+1},\ldots,\pi_n\), with \(1\leq k<n\).  Both have fewer than
\(n\) operands, so the induction hypothesis applies to each.

Suppose \((w,v)\in R_\gamma\).  By the composition clause there is a state
\(u\) with \((w,u)\in R_{\gamma_1}\) and \((u,v)\in R_{\gamma_2}\).  By the
hypothesis for \(\gamma_1\) there are states \(u_0,\ldots,u_k\) with
\(u_0=w\), \(u_k=u\) and \((u_{i-1},u_i)\in R_{\pi_i}\) for
\(1\leq i\leq k\).  By the hypothesis for \(\gamma_2\), applied to its
operands \(\pi_{k+1},\ldots,\pi_n\), there are states
\(u'_k,\ldots,u'_n\) with \(u'_k=u\), \(u'_n=v\) and
\((u'_{i-1},u'_i)\in R_{\pi_i}\) for \(k<i\leq n\).  Since \(u_k=u=u'_k\),
the sequence \(u_0,\ldots,u_k,u'_{k+1},\ldots,u'_n\) starts at \(w\),
ends at \(v\), and each consecutive pair lies in the relation of the
operand with the same index.

Conversely, let \(u_0,\ldots,u_n\) be states with \(u_0=w\), \(u_n=v\) and
\((u_{i-1},u_i)\in R_{\pi_i}\) for \(1\leq i\leq n\).  The hypothesis for
\(\gamma_1\), applied to \(u_0,\ldots,u_k\), gives
\((w,u_k)\in R_{\gamma_1}\); the hypothesis for \(\gamma_2\), applied to
\(u_k,\ldots,u_n\), gives \((u_k,v)\in R_{\gamma_2}\).  The composition
clause with the intermediate state \(u_k\) gives \((w,v)\in R_\gamma\).

\emph{Inductive step for item~2.}  Let \(n>1\) and
\(\kappa=\kappa_1\cup\kappa_2\), where \(\kappa_1\) brackets
\(\pi_1,\ldots,\pi_k\) and \(\kappa_2\) brackets
\(\pi_{k+1},\ldots,\pi_n\).  By the union clause and the induction
hypothesis,
\(R_\kappa=R_{\kappa_1}\cup R_{\kappa_2}
=(R_{\pi_1}\cup\cdots\cup R_{\pi_k})\cup(R_{\pi_{k+1}}\cup\cdots\cup R_{\pi_n})\),
and union of sets is associative, so this is
\(R_{\pi_1}\cup\cdots\cup R_{\pi_n}\).

\emph{Consequences.}  Item~1 describes \(R_\gamma\) by a condition that
mentions only \(\pi_1,\ldots,\pi_n\) and not the bracketing, so two
bracketings \(\gamma,\gamma'\) of \(\pi_1;\cdots;\pi_n\) have
\(R_\gamma=R_{\gamma'}\); item~2 gives the same for choice.  By
\cref{def:p09-pdl}, \(M,s\models[\alpha]A\) holds exactly when every
\(v\) with \((s,v)\in R_\alpha\) satisfies \(A\), a condition that depends
on \(\alpha\) only through \(R_\alpha\).  Hence
\(M,s\models[\gamma]A\) iff \(M,s\models[\gamma']A\), and likewise for two
bracketings of an iterated choice.
\end{proof}
\end{annexitem}

\begin{annexitem}{annex:p01-w02c-p01-017-prefix-chains}{thm:p01-w02c-p01-017-prefix-chains}{When the rank of a prefix operator matters}
\begin{proof}
\emph{Item~1.}  By induction on \(k\).  For \(k=1\) the string \(o_1X\) has
one operator and one operand, which is adjacent to no other operator, so
nothing competes for it and the reading is \(o_1X\).  For \(k>1\), the
operand \(X\) is adjacent only to \(o_k\), and each operator \(o_j\) with
\(j<k\) is adjacent on its right to the operator \(o_{j+1}\), not to an
operand.  A prefix operator takes the formula that follows it, and the
only formula that can follow \(o_{k-1}\) is one that begins with \(o_k\);
clause~4 of \cref{def:p01-w02-p01-024-fixity} states that two prefix
operators never compete because this nesting is forced.  So \(o_k\) takes
\(X\), and \(o_1\cdots o_{k-1}\) applies to the formula \(o_kX\), which is
an operand adjacent to no further operator on its right.  By the induction
hypothesis for the \(k-1\) operators \(o_1,\ldots,o_{k-1}\) and this
operand, the reading is \(o_1(o_2(\cdots(o_{k-1}(o_kX))\cdots))\).  No
step compares ranks, so the reading is the same for every assignment of
ranks, and when some operators have none.

\emph{Item~2.}  In \(oX\star Y\) the operand \(X\) is adjacent to \(o\)
on its left and to \(\star\) on its right, so the two operators compete
for it; \(Y\) is adjacent only to \(\star\).  The ranks differ, so
clause~3 does not apply, and clause~1 gives \(X\) to the operator with the
smaller rank.  If \(r<s\), the prefix operator takes \(X\), forming
\(oX\); this formula is then the left operand of \(\star\), adjacent to no
other operator, and the reading is \((oX)\star Y\).  If \(r>s\), the
binary operator takes \(X\), forming \(X\star Y\); the prefix operator
then takes the shortest following formula, which is now \(X\star Y\)
because \(X\) alone is already an operand of \(\star\), and the reading
is \(o(X\star Y)\).  When \(o\) has no rank, clause~1 cannot compare the
competitors, clause~3 concerns equal ranks, and clause~4 concerns two
unary operators; no clause decides the competition, so the string has no
licensed reading and parentheses are needed.
\end{proof}
\end{annexitem}

\begin{annexitem}{annex:p01-w02c-p01-017-printed-readings}{thm:p01-w02c-p01-017-printed-readings}{Readings of the printed Part~9 and Part~19 formulas}
\begin{proof}
\emph{Item~1.}  In \([a;p?]q\) the program between the brackets is
\(a;p?\).  By item~1 of \cref{def:p01-w02c-p01-017-readings} the operand
of \(?\) is the longest formula ending before it; the candidates are
\(p\) and \(a;p\), and \(a;p\) is not a formula because \(a\) is a
program, so the operand is \(p\) and the test is \(p?\).  The remaining
program is the composition of \(a\) and \(p?\), so the formula is
\([a;(p?)]q\).  In \([p?;c]q\) the only formula before \(?\) is \(p\),
giving \([(p?);c]q\).  In \([a;b^{*}]q\) item~2 applies the star to the
immediately preceding atomic program \(b\), and composition joins \(a\)
and \(b^{*}\), giving \([a;(b^{*})]q\).  In each case the brackets close
the program, and the box applies to the following letter \(q\).

\emph{Item~2.}  The formula operators involved are the box (rank~3,
prefix, its program delimited by its brackets), conjunction (rank~4, left
associative), implication (rank~6, right associative) and the
biconditional (rank~7, nonassociative), by the registry row of the box
and \cref{def:p01-w02-p01-026-precedence}.

In \([\alpha^{*}]A\leftrightarrow(A\land[\alpha][\alpha^{*}]A)\) the
letter \(A\) after the first box is adjacent to that box and to
\(\leftrightarrow\); by clause~1 of \cref{def:p01-w02-p01-024-fixity}
the box (3) takes it before the biconditional (7).  Inside the
parentheses, \([\alpha][\alpha^{*}]A\) is a chain of two prefix
operators, which nests by item~1 of
\cref{thm:p01-w02c-p01-017-prefix-chains} to
\([\alpha]([\alpha^{*}]A)\); the first letter \(A\) inside the
parentheses is adjacent only to \(\land\), whose right operand is the
nested box formula.  The whole is
\(([\alpha^{*}]A)\leftrightarrow(A\land([\alpha]([\alpha^{*}]A)))\).

In \([\alpha^{*}](A\to[\alpha]A)\to(A\to[\alpha^{*}]A)\) the box
\([\alpha^{*}]\) is followed by a parenthesized formula, which is one
operand; it is adjacent to the box and to the main \(\to\), and the box
(3) takes it before the implication (6).  Inside the first parentheses,
the last \(A\) is taken by \([\alpha]\) (3) rather than by \(\to\) (6),
giving \(A\to([\alpha]A)\); the same holds in the second parentheses.
The reading is \(([\alpha^{*}](A\to([\alpha]A)))\to(A\to([\alpha^{*}]A))\).

In \([B?]A\leftrightarrow(B\to A)\) the program \(B?\) is delimited by the
brackets, the box takes \(A\) before the biconditional by clause~1, and
the reading is \(([B?]A)\leftrightarrow(B\to A)\).

In \(X\land[\alpha^{*}](X\to[\alpha]X)\to[\alpha^{*}]X\) the box
\([\alpha^{*}]\) takes the parenthesized formula, adjacent to it and to
the main \(\to\), by clause~1 (3 against 6), and inside the parentheses
\([\alpha]\) takes \(X\) as before.  The box formula so formed is adjacent
to \(\land\) and to \(\to\); by clause~1 conjunction (4) takes it before
implication (6), giving \(X\land[\alpha^{*}](X\to([\alpha]X))\).  The last
\(X\) is taken by the final box (3) before \(\to\) (6).  The reading is
\((X\land([\alpha^{*}](X\to([\alpha]X))))\to([\alpha^{*}]X)\).

\emph{Item~3.}  In
\(A_1\land\cdots\land A_m\land\neg LB_1\land\cdots\land\neg LB_n\) each
\(\neg LB_j\) is a chain of two prefix operators before the letter
\(B_j\).  The letter is adjacent to \(L\) on its left and, unless \(j=n\),
to \(\land\) on its right.  By item~6 of
\cref{def:p01-w02c-p01-017-readings}, \(L\) has rank~3, below the rank~4
of conjunction, so by item~2 of
\cref{thm:p01-w02c-p01-017-prefix-chains} \(L\) takes \(B_j\); by item~1
of the same lemma the negation applies to \(LB_j\), giving
\(\neg(LB_j)\), and this formula is adjacent only to conjunction signs.
Each \(A_i\) is a letter adjacent only to conjunction signs.  The string
is therefore an iterated conjunction of the \(m+n\) conjuncts in the
order displayed, and clause~2 of \cref{def:p01-w02-p01-024-fixity} with
the left association of \(\land\) gives the left-associated bracketing.
In \(LF\) and \(\neg LF\) the operand of \(L\) is a formula metavariable,
which denotes one formula and is therefore a delimited operand; by item~1
of \cref{thm:p01-w02c-p01-017-prefix-chains}, \(LF\) is \(L\) applied to
\(F\) and \(\neg LF\) is \(\neg(LF)\).

\emph{Item~4.}  The denials \(\neg q\) and \(\neg p\) are strings
\(oX\) with one prefix operator and a letter.  By item~1 of
\cref{thm:p01-w02c-p01-017-prefix-chains} with \(k=1\), the reading is
\(\neg q\), respectively \(\neg p\), whatever the rank of \(\neg\), so the
ranks 60 and~3 agree.  For \(\neg X\land Y\), item~2 of the same lemma
with \(s=4\) gives \(\neg(X\land Y)\) when \(r=60>4\) and
\((\neg X)\land Y\) when \(r=3<4\).  These are different formulas, and
they are not equivalent: for letters \(X,Y\) under the valuation that
makes \(X\) true and \(Y\) false, \(\neg(X\land Y)\) is true and
\((\neg X)\land Y\) is false.

\emph{Item~5.}  By item~8 of \cref{def:p01-w02c-p01-017-readings} the
diamond is a prefix operator of rank~3 whose program is closed by the
angle brackets, so each \(\langle\pi\rangle\) acts on formulas like a
single prefix sign.  The other operators are the global negation
(rank~3, prefix), conjunction (rank~4), disjunction (rank~5),
implication (rank~6) and the biconditional (rank~7), and the box of
rank~3.  Two facts decide every string.  (a)~By item~1 of
\cref{thm:p01-w02c-p01-017-prefix-chains}, a chain of prefix signs
before an operand nests, whatever the ranks.  (b)~By item~2 of the same
lemma, an operand adjacent to a prefix sign of rank~3 on its left and to
a binary connective of rank at least~4 on its right is taken by the
prefix sign.  A parenthesized formula is a single operand.

In \(\neg\langle\alpha\rangle\neg X\) the letter \(X\) follows the chain
\(\neg,\langle\alpha\rangle,\neg\) and nothing follows it, so (a) gives
\(\neg(\langle\alpha\rangle(\neg X))\).

Axiom (2): \(\bot\) after \(\langle\alpha\rangle\) is adjacent to
\(\leftrightarrow\) (7) on its right, so (b) gives
\(\langle\alpha\rangle\bot\) as the left operand of the biconditional,
whose right operand is \(\bot\).

Axiom (3): the program \(\alpha\cup\beta\) is enclosed by the angle
brackets; \(X\) is taken by the diamond before \(\leftrightarrow\) by
(b).  Inside the parentheses on the right, the first \(X\) is between
\(\langle\alpha\rangle\) and \(\lor\) (5) and goes to the diamond by
(b), and the last \(X\) follows \(\langle\beta\rangle\) and ends the
parenthesized formula, so it goes to that diamond.  The parenthesized
formula is the disjunction of the two diamond formulas.

Axiom (4): the left side is a diamond applied to the parenthesized
formula \(X\lor Y\), which is one operand adjacent to
\(\leftrightarrow\) and taken by the diamond by (b); the right side is
read as in axiom (3).

Axiom (5): on the left, \(X\) is taken by \(\langle\alpha;\beta\rangle\)
before \(\leftrightarrow\) by (b).  On the right,
\(\langle\alpha\rangle\langle\beta\rangle X\) is a chain of two prefix
signs before the last letter, which (a) nests as
\(\langle\alpha\rangle(\langle\beta\rangle X)\).

Axiom (6): the left side is read as in axiom (5).  Inside the
parentheses, the first \(X\) is a letter adjacent only to \(\lor\), and
\(\langle\alpha\rangle\langle\alpha^{*}\rangle X\) nests by (a); so the
right side is \(X\lor(\langle\alpha\rangle(\langle\alpha^{*}\rangle X))\).

Axiom (7): \(X\) after \(\langle\alpha^{*}\rangle\) is adjacent to
\(\to\) (6) and goes to the diamond by (b).  The next \(X\) is adjacent
to \(\to\) on its left and to \(\lor\) (5) on its right; by clause~1 of
\cref{def:p01-w02-p01-024-fixity} disjunction takes it, so the
consequent of the implication is a disjunction whose left operand is
\(X\).  The right operand of the disjunction is
\(\langle\alpha^{*}\rangle\) applied to the parenthesized formula, which
ends the string, so the diamond takes it.  Inside the parentheses,
\(X\) after \(\neg\) is adjacent to \(\land\) (4) and goes to the
negation by (b), and the last \(X\) follows \(\langle\alpha\rangle\) and
ends the parenthesized formula; so it is \((\neg X)\land(\langle\alpha\rangle X)\).
Since the implication (6) has a larger rank than the disjunction (5),
the disjunction is formed first and is the whole consequent, which gives
the displayed reading.

For the remaining formulas: in \(\langle a\rangle Q\lor\langle b\rangle Q\)
the first \(Q\) is taken by \(\langle a\rangle\) before \(\lor\) by (b),
and the second ends the string after \(\langle b\rangle\).  The formula
\(\neg\langle a\rangle\neg Q\) is read as the definition above.  In
\(C\land\langle a\rangle D\) the letter \(C\) is adjacent only to
\(\land\), and \(D\) ends the string after \(\langle a\rangle\).  In
\(Y\land\langle\beta\rangle\neg Y\) the last \(Y\) follows the chain
\(\langle\beta\rangle,\neg\), which nests by (a), and the first \(Y\) is
adjacent only to \(\land\).  In \(\langle a\rangle A\to[a]A\) the first
\(A\) is taken by \(\langle a\rangle\) before \(\to\) by (b), and the
second ends the string after the box.

\emph{Item~6.}  By item~9 of \cref{def:p01-w02c-p01-017-readings} the
diamond and the box of the modal \(\mu\)-calculus are prefix signs of
rank~3.  In \(\langle a\rangle p\land[a]q\) the letter \(p\) is between
\(\langle a\rangle\) and \(\land\) (4) and goes to the diamond by (b),
and \(q\) ends the string after the box, so the formula is
\((\langle a\rangle p)\land([a]q)\).  In
\(\mu Z.(\langle a\rangle Y\lor[b]Z)\) the binder applies to the
parenthesized formula, which ends the string; inside it, \(Y\) is
between \(\langle a\rangle\) and \(\lor\) (5) and goes to the diamond by
(b), and \(Z\) ends the parenthesized formula after the box.  The formula
\(\neg\langle a\rangle\neg p\) is a chain of three prefix signs before
\(p\), which nests by (a).
\end{proof}
\end{annexitem}

\section{Proofs for the readings of Part 8 operators in their presenting chapters}
\label{sec:annex-p01-w02c-p01-017-p8-chapters}

\begin{annexitem}{annex:p01-w02c-p01-017-chapter-readings}{thm:p01-w02c-p01-017-chapter-readings}{Readings of printed formulas of the Part~8 presenting chapters}
\begin{proof}
The ranks used are those of \cref{def:p01-w02c-p01-017-p8-chapters}:
every prefix operator that occurs has rank~3, until has rank~3.5,
\({}^\perp\) has rank~2.5, \(\land\), \(\otimes\) and \(\bindnasrepma\)
have rank~4, \(\lor\) has rank~5, \(\to\) and \(\multimap\) have rank~6
and \(\leftrightarrow\) has rank~7.  Four facts are used throughout.
\begin{enumerate}
\item[(a)] A chain of prefix operators before an operand nests, whatever
  the ranks (item~1 of \cref{thm:p01-w02c-p01-017-prefix-chains}).
\item[(b)] If a chain \(o_1\cdots o_k\) of prefix operators, each of
  rank smaller than the rank \(s\) of a binary operator \(\star\), stands
  before an operand \(X\) that is followed by \(\star\), the string
  \(o_1\cdots o_kX\star Y\) is \((o_1(\cdots(o_kX)\cdots))\star Y\).  For
  \(k=1\) this is item~2 of \cref{thm:p01-w02c-p01-017-prefix-chains}.
  For \(k>1\), the operand \(X\) is adjacent to \(o_k\) and to \(\star\),
  and \(o_k\) takes it by clause~1 of \cref{def:p01-w02-p01-024-fixity};
  the formula \(o_kX\) is then an operand between \(o_{k-1}\) and
  \(\star\), and the claim follows by induction on \(k\).
\item[(c)] A parenthesized formula is one operand, and an operand that
  ends the string, or ends the parenthesized formula that contains it,
  is adjacent only to the operator on its left.
\item[(d)] Of two binary operators of different ranks that compete for
  an operand, the one with the smaller rank takes it (clause~1 of
  \cref{def:p01-w02-p01-024-fixity}).
\end{enumerate}

\emph{Item~1.}  In \(\mathsf X(A\to B)\to(\mathsf XA\to\mathsf XB)\), the
parenthesized \((A\to B)\) is adjacent to \(\mathsf X\) and to the main
\(\to\); next (3) takes it by (b).  Inside the second parentheses, \(A\)
lies between \(\mathsf X\) (3) and \(\to\) (6) and goes to \(\mathsf X\)
by (b), and \(B\) ends the parentheses after \(\mathsf X\) by (c).  The
reading is \((\mathsf X(A\to B))\to((\mathsf XA)\to(\mathsf XB))\).  In
\(\mathsf X\neg A\leftrightarrow\neg\mathsf XA\), the chain
\(\mathsf X,\neg\) before \(A\) and the biconditional give
\(\mathsf X(\neg A)\) by (a) and (b), and the chain \(\neg,\mathsf X\)
before the final \(A\) nests by (a) and (c), giving
\(\neg(\mathsf XA)\).

In \(A\mathbin{\mathsf U}B\leftrightarrow B\lor(A\land\mathsf
X(A\mathbin{\mathsf U}B))\), the first \(A\) is adjacent only to until.
The first \(B\) is adjacent to until (3.5) and to \(\leftrightarrow\)
(7), and until takes it by (d), forming \(A\mathbin{\mathsf U}B\).  The
second \(B\) is adjacent to \(\leftrightarrow\) (7) and to \(\lor\) (5),
and \(\lor\) takes it by (d).  The right operand of \(\lor\) is the
parenthesized formula, which ends the string, by (c).  Inside it, \(A\) is
adjacent only to \(\land\), and \(\mathsf X\) is followed by the
parenthesized \((A\mathbin{\mathsf U}B)\), which ends the enclosing
parentheses, so \(\mathsf X\) takes it by (c) and the formula
\(\mathsf X(A\mathbin{\mathsf U}B)\) is the right operand of \(\land\).
The biconditional joins the two sides, which gives
\((A\mathbin{\mathsf U}B)\leftrightarrow(B\lor(A\land(\mathsf
X(A\mathbin{\mathsf U}B))))\).

In the premise \(B\lor(A\land\mathsf XC)\to C\) of the induction rule,
\(B\) is adjacent only to \(\lor\); the parenthesized formula is adjacent
to \(\lor\) (5) and to \(\to\) (6), and \(\lor\) takes it by (d).  Inside
it \(C\) ends the parentheses after \(\mathsf X\), so the inner formula is
\(A\land(\mathsf XC)\).  The final \(C\) is adjacent only to \(\to\).  In
the conclusion \(A\mathbin{\mathsf U}B\to C\), \(B\) lies between until
(3.5) and \(\to\) (6) and goes to until by (d).  The definiens
\(\topv\mathbin{\mathsf U}A\) has one operator, and
\(\neg\mathsf F\neg A\) is a chain of three prefix operators before
\(A\), which nests by (a) to \(\neg(\mathsf F(\neg A))\).

\emph{Item~2.}  The first schema has the shape of the first schema of
item~1 with \(K_i\) in place of \(\mathsf X\), and the same steps give its
reading.  In \(K_iA\to A\) the first \(A\) lies between \(K_i\) (3) and
\(\to\) (6) and goes to \(K_i\) by (b).  In \(K_iA\to K_iK_iA\) the left
side is read in the same way, and the right side is a chain of two prefix
operators that ends the string, which nests by (a).  In
\(\neg K_iA\to K_i\neg K_iA\) the chain \(\neg,K_i\) before \(A\) is
followed by \(\to\), so (b) gives \(\neg(K_iA)\) as the antecedent; the
consequent is the chain \(K_i,\neg,K_i\) before the final \(A\), which
nests by (a) to \(K_i(\neg(K_iA))\).  In
\(C_GA\to E_G(A\land C_GA)\), the first \(A\) goes to \(C_G\) by (b); the
parenthesized formula follows \(E_G\) and ends the string, so \(E_G\)
takes it by (c); inside it, \(A\) is adjacent only to \(\land\), and the
last \(A\) ends the parentheses after \(C_G\).  In \(p\land\neg K_ip\) the
letter \(p\) is adjacent only to \(\land\), and the chain \(\neg,K_i\)
before the final \(p\) ends the string and nests by (a); the nested
formula is the right operand of \(\land\).

\emph{Item~3.}  The definiens \(\neg\mathsf A\neg\alpha\) is a chain of
three prefix operators before \(\alpha\), and \(\mathsf G\mathsf Fp\) a
chain of two before \(p\); both nest by (a).

\emph{Item~4.}  In \(\Diamond\Diamond p\to\Box(\Diamond p\lor\Box p)\) the
chain \(\Diamond,\Diamond\) before \(p\) is followed by \(\to\), and (b)
gives \(\Diamond(\Diamond p)\).  The parenthesized formula follows
\(\Box\) and ends the string, so \(\Box\) takes it by (c).  Inside it,
\(p\) lies between \(\Diamond\) (3) and \(\lor\) (5) and goes to the
diamond by (b), and the last \(p\) ends the parentheses after \(\Box\).
In \(\Box p\lor\Box\neg p\) the first \(p\) goes to \(\Box\) by (b), and
the chain \(\Box,\neg\) before the last \(p\) nests by (a).

\emph{Item~5.}  In \(A^\perp\bindnasrepma B^\perp\) the letter \(A\) is
adjacent only to the postfix sign that follows it, which takes it, and
the formula \(A^\perp\) is the left operand of par.  The letter \(B\) is
adjacent to par (4) on its left and to the postfix \({}^\perp\) (2.5) on
its right; by clause~1 of \cref{def:p01-w02-p01-024-fixity} the postfix
sign takes it, so the reading is \((A^\perp)\bindnasrepma(B^\perp)\).  The
same steps read \(A^\perp\bindnasrepma B\) as \((A^\perp)\bindnasrepma
B\), the final \(B\) being adjacent only to par.  In
\(A\multimap B\multimap C\) both competitors for \(B\) are occurrences of
\(\multimap\), whose association is right, so clause~2 of
\cref{def:p01-w02-p01-024-fixity} gives \(A\multimap(B\multimap C)\).  In
\(A\otimes B\multimap C\) the letter \(B\) lies between \(\otimes\) (4)
and \(\multimap\) (6) and goes to the tensor by (d).
\end{proof}
\end{annexitem}

\begin{annexitem}{annex:p01-w02c-p01-017-cll-negation}{thm:p01-w02c-p01-017-cll-negation}{Only the source reading of the exponential clauses defines linear negation}
\begin{proof}
\emph{Item~1.}  Every formula is an atom, a unit, or a compound with
exactly one outermost constructor and its immediate subformulas, so
definitions by recursion on formulas are available: a value may be given
outright at atoms and units, and at a compound formula as a function of
the values at its immediate subformulas.  Under the reading of item~10 of
\cref{def:p01-w02c-p01-017-p8-chapters}, each right-hand side applies
\({}^\perp\) only to an immediate subformula \(A\) or \(B\) of the
left-hand side and then applies one constructor: \(\bindnasrepma\),
\(\otimes\), \(\oplus\), \(\mathbin{\&}\), \(?\) or \(!\).  The clauses
therefore have the form of such a recursion, and they define an operation
\(f\) on all formulas.  If \(g\) is any operation satisfying the same
equations, then \(g(A)=f(A)\) for every formula \(A\), by induction on
\(A\): at atoms and units both take the listed value, and at a compound
formula both apply the same constructor to values at the immediate
subformulas, which agree by the induction hypothesis.  So the operation is
unique.

We show \(f(f(A))=A\) by induction on \(A\).  For atoms,
\(f(f(p))=f(p^\perp)=p\) and \(f(f(p^\perp))=f(p)=p^\perp\).  For units,
\(f(f(\mathbf 1))=f(\bot)=\mathbf 1\), \(f(f(\bot))=f(\mathbf 1)=\bot\),
\(f(f(\top))=f(0)=\top\) and \(f(f(0))=f(\top)=0\).  For a tensor,
\[
 f(f(A\otimes B))=f(f(A)\bindnasrepma f(B))=f(f(A))\otimes f(f(B))
 =A\otimes B,
\]
by the tensor clause, the par clause and the induction hypothesis for
\(A\) and \(B\).  The par, with and plus cases are the same computation
with the clauses of the dual pair.  For the exponentials,
\[
 f(f(!A))=f(?f(A))=!f(f(A))=!A,
 \qquad
 f(f(?A))=f(!f(A))=?f(f(A))=?A,
\]
by the bang clause, the why-not clause and the induction hypothesis.

\emph{Item~2.}  Under the alternative reading the two exponential clauses
say the same thing, \(g(!A)=g(?A)\), for an operation \(g\) on formulas.
An involution is injective: if \(g(x)=g(y)\) then
\(x=g(g(x))=g(g(y))=y\).  So an involution satisfying the clause for some
\(A\) would give \(!A=?A\); but \(!A\) and \(?A\) have different outermost
constructors, so they are different formulas.  Hence no involution
satisfies the clauses.

For the second claim fix a formula \(A\) and two different formulas
\(C\) and \(C'\).  Define \(g_C\) by recursion on formulas: at atoms,
units and binary compounds by the clauses of item~1, at the two formulas
\(!B\) and \(?B\) by the value \(C\) when \(B\) is \(A\), and by the fixed
formula \(\mathbf 1\) when \(B\) is another formula.  The value at an
exponential formula does not depend on other values, so this is a
recursion of the kind used in item~1, and \(g_C\) satisfies every clause:
the atom, unit and binary clauses by construction, and both exponential
clauses because \(g_C(!B)=g_C(?B)\) for every \(B\).  Define \(g_{C'}\) in
the same way with \(C'\).  Both operations satisfy all the equations, and
\(g_C(!A)=C\), which differs from \(C'=g_{C'}(!A)\).  So the equations do
not determine the value at \(!A\).
\end{proof}
\end{annexitem}
